\section{Background}
\label{sec:background}

\subsection{ReCom}

The Recombination (\recom{}) Markov chain~\citep{deford2021} is the
standard Markov chain Monte Carlo method for redistricting.
Given a current plan $\pi$ on census-tract graph $G = (V, E)$:
\begin{sloppypar}
\begin{enumerate}
\item Select a random boundary edge $(u, v)$ with $\pi(u) \neq \pi(v)$.
\item Merge districts $A = \pi^{-1}(\pi(u))$ and $B = \pi^{-1}(\pi(v))$
      into the subgraph $G[A \cup B]$.
\item Sample a random spanning tree $\mathcal{T}$ of $G[A \cup B]$.
\item Find an edge in $\mathcal{T}$ whose removal partitions $A \cup B$
      into two population-balanced subsets; if none exists, reject.
\item Update $\pi$ with the new bipartition.
\end{enumerate}
\end{sloppypar}

\recom{} is designed to have a stationary distribution approximately
uniform over valid plans (connected, population-balanced bipartitions).
It is the foundational chain for all methods discussed in this paper.

\subsection{ConvergenceSweep}

\cs{}~\citep{b16paper} runs a single \recom{} chain for $T$ steps
and returns the best plan seen during the run.
The parameter $T = 600$ is calibrated as the statutory stopping threshold:
chains run beyond $T = 600$ steps show diminishing Polsby-Popper improvement
(less than 1\% per additional 100 steps) on all 50 states.

\cs{} operates with a fixed step budget and selects the minimum-edge-cut
plan encountered during the entire run.
Selecting the within-run minimum is appropriate for \cs{} because the
chain is run as a single continuous trajectory---there is no restart
mechanism, and the chain is not being used as a Markov restart point.

\subsection{BisectionEnsemble}

\be{}~\citep{h1paper} applies \recom{} ensembles at each node of the
\ar{} bisection tree.
Rather than running a single ensemble on the full state, \be{} runs
a short ensemble ($T = 100$ steps) at each bisection node, selecting
the best endpoint at each node.
The node-level best plan propagates down the tree.

\be{} with $p = 0.0$ (always select the best endpoint) achieves
15--26\% Polsby-Popper improvement over single-seed METIS across
NC, WI, and TX, and is the strongest benchmark in our comparison.

\subsection{Short-Burst: Original Paper}

Cannon, Duchin, Randall, and Rule~\citep{cannon2022} introduced
Short-Burst optimization in the context of partisan outcome analysis.
Their setting: given an enacted redistricting plan, measure its partisan
properties by running short \recom{} chains from the enacted plan
and examining the distribution of partisan outcomes at chain endpoints.

Short bursts of length $\ell$ starting from a given plan sample the
\emph{local neighbourhood} of that plan in plan space.
By restarting from each endpoint rather than the initial plan, successive
bursts explore the neighbourhood of the neighbourhood, gradually
drifting toward better plans without the slow mixing of a long chain.

\paragraph{Endpoint vs.\ minimum.}

A critical design choice distinguishes Short-Burst from naive greedy
chain-and-pick:
Cannon et al.\ retain the \emph{endpoint} of each burst, not the
within-burst minimum.

This distinction has two consequences:

\begin{enumerate}
\item \textbf{Markov property.}
      The endpoint of a \recom{} chain is a valid plan in $\mathcal{P}$
      that could be the starting point of a new chain with the correct
      Markov structure.
      The within-burst minimum is not a valid restart point in general:
      it was not the chain's terminal state, and the chain did not end
      in that state---restarting from it would break the Markov property
      and introduce path dependencies not accounted for in the stationary
      distribution.

\item \textbf{Stationarity.}
      Selecting the minimum within each burst would systematically bias
      the restarted chain toward lower-edge-cut plans in a way not captured
      by the chain's stationary distribution.
      The bias is a form of selection pressure: the chain preferentially
      transitions to lower-cut plans not because the transition probability
      is higher, but because the selection mechanism discards higher-cut endpoints.
      This makes the selected sequence a biased chain, not a valid
      Markov chain on $\mathcal{P}$.
\end{enumerate}

In our min-EC compactness setting, the same argument applies: we retain
burst endpoints and select among them using the percentile parameter $p$~\citep{h0paper}.

\paragraph{From partisan analysis to compactness optimisation.}

Cannon et al.\ used Short-Burst to measure partisan sensitivity: how
much does a plan's partisan lean vary in its local neighbourhood?
Our use is different: we use Short-Burst as an optimization algorithm
to find low-edge-cut plans.
The mechanism is identical; only the objective differs.
This is a natural extension---the endpoint-retention and restart properties
are valuable regardless of the objective being optimized.
